\section{The axial gauge and the gauge-adapted averaging}\label{sec:adapted}

The linear theory of Parts~II and~III is formulated in Ba\l aban's Landau gauge, with the residual gauge transformations $\cN(Q'(U))$, i.e.\ those with vanishing \emph{averages}. The nonlinear averaging of Section~\ref{sec:averaging} is covariant with respect to the \emph{values} of gauge transformations at the points of $\T$. Ba\l aban bridges the two by an intermediate gauge, the axial gauge, and by a modified averaging \cite[Sect.~C]{BalAvg}. This section develops that construction in a form adapted to our needs. The main results are the correspondence theorem (Theorem~\ref{thm:correspondence}), the expansion of the gauge-adapted averaging (Proposition~\ref{prop:Q-expansion}), and the properties of its linearization (Proposition~\ref{prop:linearized-averages}). The last one shows that Parts~II and~III apply to Ba\l aban's averaging.

\subsection{The reference configuration and relative configurations}

Fix a real configuration $U_0\in\mathfrak U_0(a)$ on $\Tf$ with $a\le c_1$, and put $W_j=U_0^{(j)}$. By Proposition~\ref{prop:avg-regular}, $W_j\in\mathfrak U_j(2a)$. A \emph{relative configuration} at scale $j$ is a configuration $X$ on $\T^{(j)}$ close to $1$; it represents the configuration $XW_j$ (bondwise product). We use the following notation.
\begin{itemize}
\item \emph{Relative transports.} $X_W(\Gamma)=(XW_j)(\Gamma)\,W_j(\Gamma)^{-1}$. For $\Gamma=b_1\cdots b_m$ this is $\prod_iR(W_j(\Gamma_{<i}))X(b_i)$, with $\Gamma_{<i}=b_1\cdots b_{i-1}$ and $R(Y)Z=YZY^{-1}$.
\item \emph{Relative gauge action.} $X^{(g)}=(XW_j)^gW_j^{-1}$, i.e.\ $X^{(g)}(b)=g(b_-)X(b)\,R(W_j(b))\big(g(b_+)^{-1}\big)$. Then $(X^{(g)})^{(h)}=X^{(hg)}$, and $X^{(g)}_W(\Gamma)=g(\Gamma_-)\,X_W(\Gamma)\,R(W_j(\Gamma))\big(g(\Gamma_+)^{-1}\big)$. The configurations $1^{(g)}$ are the \emph{pure gauges}.
\item \emph{Relative averages.} For $X$ at scale $j-1$ put $\overline X=\mathfrak a_j(XW_{j-1})W_j^{-1}$. By Lemma~\ref{lem:one-step-basic}(b), $\overline{X^{(g)}}=\overline X^{(g|_{\T^{(j)}})}$ for real $g$. For a fine relative configuration $U'$ put $U'^{[j]}=(U'U_0)^{(j)}W_j^{-1}$, the iterated relative average.
\item \emph{Lie fields.} We write $X=e^{\ii\delta_jB}$, and $D_{W_j}\lambda(b)=\delta_j^{-1}\big(R(W_j(b))\lambda(b_+)-\lambda(b_-)\big)$. At $j=0$ this is the covariant derivative $D_{U_0}$ of Part~II; at $j=k$ it is the unit-lattice derivative $D_V$ with $V=W_k$. To first order, $(e^{\ii\delta_jB})^{(e^{\ii\lambda})}=e^{\ii\delta_j(B-D_{W_j}\lambda)}$. On reversed bonds the relative configuration is $X(-b)=R(W_j(b))^{-1}(X(b)^{-1})$, since $(XW_j)(-b)=(XW_j)(b)^{-1}$; correspondingly $B(-b)=-R(W_j(b))^{-1}B(b)$. Only positively oriented bonds carry independent variables.
\end{itemize}

\subsection{The axial gauge}

\begin{definition}[Axial gauge {\cite[(1.19)]{BalReg}}]\label{def:axial}
A relative configuration $X$ at scale $j-1$ is \emph{axial at scale $j$} if $X_W(\Gamma^j_{y,x})=1$ for all $y\in\T^{(j)}$ and $x\in B_j(y)$. A fine relative configuration $U'$ belongs to $\mathrm{Ax}_j$ if $U'^{[i-1]}$ is axial at scale $i$ for $i=1,\dots,j$.
\end{definition}

For $h$ on $\T^{(j)}$ let $\mathcal E_j(h)$ be the function on $\T^{(j-1)}$ given by $\mathcal E_j(h)(x)=R(W_{j-1}(\Gamma^j_{y,x}))^{-1}h(y)$ for $x\in B_j(y)$, and put $\mathcal E^{(j)}=\mathcal E_1\circ\cdots\circ\mathcal E_j$. This is the extension of $h$ from $\T^{(j)}$ to $\Tf$ by parallel transport with the background along the trees. It is multiplicative, $\mathcal E^{(j)}(hh')=\mathcal E^{(j)}(h)\mathcal E^{(j)}(h')$, and it does not change values at the points of $\T^{(j)}$.

\begin{lemma}[Axial transformations]\label{lem:axial}
\begin{enumerate}[label=(\alph*)]
\item Let $X$ be a relative configuration at scale $j-1$ and $h$ a function on $\T^{(j)}$. There is exactly one $g$ on $\T^{(j-1)}$ with $g|_{\T^{(j)}}=h$ and $X^{(g)}$ axial at scale $j$, namely
\[
g(x)=R(W_{j-1}(\Gamma^j_{y,x}))^{-1}\big(h(y)\,X_W(\Gamma^j_{y,x})\big),\qquad x\in B_j(y).
\]
If $X$ is axial at scale $j$, this is $g=\mathcal E_j(h)$.
\item Let $U'$ be a real fine relative configuration whose iterated relative averages $U'^{[i]}$, $i<j$, are defined, and $h$ a real function on $\T^{(j)}$. There is exactly one real $g$ on $\Tf$ with $g|_{\T^{(j)}}=h$ and $U'^{(g)}\in\mathrm{Ax}_j$. If $U'\in\mathrm{Ax}_j$, this is $g=\mathcal E^{(j)}(h)$.
\item Every orbit of $\mathfrak G_0$ in $\mathfrak U_0(c_1)$ contains exactly one $U$ with $UU_0^{-1}\in\mathrm{Ax}_k$.
\end{enumerate}
\end{lemma}

\begin{proof}
(a) By the transformation law of relative transports, $X^{(g)}_W(\Gamma^j_{y,x})=1$ iff
\[
R(W_{j-1}(\Gamma^j_{y,x}))g(x)=g(y)X_W(\Gamma^j_{y,x}).
\] (b) Let $g$ be real. Conjugation by $G$-valued elements preserves the domain of $\ell$, so the relative averages of $U'^{(g)}$ are defined and covariant, $(U'^{(g)})^{[i-1]}=(U'^{[i-1]})^{(g|_{\T^{(i-1)}})}$ (Lemma~\ref{lem:one-step-basic}(b)). By (a), its axiality at scale $i$ determines $g|_{\T^{(i-1)}}$ uniquely from $g|_{\T^{(i)}}$, and the formula of (a) shows that it is real. Proceed from $i=j$ down to $i=1$. (c) is (b) with $j=k$, $h=1$ and $U'=UU_0^{-1}$, whose iterated relative averages are defined by Proposition~\ref{prop:avg-regular}.
\end{proof}

\subsection{The gauge-adapted averaging and the correspondence theorem}

For a relative configuration $X$ at scale $j-1$ let
\[
\tau_X(y)=\exp\Big(\ii L^{-d}\sum_{x\in B_j(y)}\ell\big(X_W(\Gamma^j_{y,x})\big)\Big),\qquad y\in\T^{(j)},
\]
be the averaged contour variable. Let $g_X$ be the transformation of Lemma~\ref{lem:axial}(a) with $h=\tau_X^{-1}$.

\begin{definition}[Gauge-adapted averaging]\label{def:adapted}
The \emph{modified one-step average} is $\tilde{\mathfrak a}_j(X)=\overline{X^{(g_X)}}=\overline X^{(\tau_X^{-1})}$, i.e.
\[
\tilde{\mathfrak a}_j(X)(c)=\tau_X(c_-)^{-1}\,\overline X(c)\,R(W_j(c))\big(\tau_X(c_+)\big).
\]
The \emph{gauge-adapted averaging} is $\mathcal Q_0=\mathrm{id}$, $\mathcal Q_j=\tilde{\mathfrak a}_j\circ\mathcal Q_{j-1}$, and $\mathcal Q=\mathcal Q_k$. For a fine relative configuration $U_1$ put $X_i=\mathcal Q_i(U_1)$ and define the \emph{correspondence transformation}
\[
\mathsf w_0(U_1)=1,\qquad\mathsf w_j(U_1)=\mathcal E^{(j-1)}\big(g_{X_{j-1}}\big)\,\mathsf w_{j-1}(U_1).
\]
Finally let $\mathfrak m_j(g)(y)=g(y)\exp\big(\ii L^{-d}\sum_{x\in B_j(y)}\ell(g(y)^{-1}R(W_{j-1}(\Gamma^j_{y,x}))g(x))\big)$ and $\mathfrak M^{(j)}=\mathfrak m_j\circ\cdots\circ\mathfrak m_1$; these are Ba\l aban's averages of gauge transformations \cite[(78)--(80)]{BalAvg}.
\end{definition}

This is Ba\l aban's construction \cite[(58)--(63), (88)--(92)]{BalAvg}: the modified average first moves $X$ to the axial gauge at scale $j$, with a normalization chosen at the coarse points, and then averages. All maps are defined when the configurations involved are sufficiently close to $1$; precise domains are given in Lemma~\ref{lem:one-step-expansion}.

\begin{theorem}[Correspondence]\label{thm:correspondence}
Let $U_1$ be a real fine relative configuration for which the maps above are defined. In~(b), (e) and~(f) the functions $h$ and $g$ are real.
\begin{enumerate}[label=(\alph*)]
\item $U_1^{(\mathsf w_j(U_1))}\in\mathrm{Ax}_j$ and $\big(U_1^{(\mathsf w_j(U_1))}\big)^{[j]}=\mathcal Q_j(U_1)$.
\item The transformations $g$ with $U_1^{(g)}\in\mathrm{Ax}_j$ are exactly $g=\mathcal E^{(j)}(h)\,\mathsf w_j(U_1)$, with $h$ a function on $\T^{(j)}$.
\item At the points $y\in\T^{(j)}$: $\mathsf w_j(U_1)(y)=\tau_{X_{j-1}}(y)^{-1}\tau_{X_{j-2}}(y)^{-1}\cdots\tau_{X_0}(y)^{-1}$.
\item If $X_{j-1}$ is axial at scale $j$, then $\mathsf w_j(U_1)=\mathsf w_{j-1}(U_1)$ and $\mathcal Q_j(U_1)=\overline{X_{j-1}}$.
\item \emph{(Pure gauges.)} If the maps are defined at $1^{(g)}$, then $\mathcal Q_j(1^{(g)})=1^{(\mathfrak M^{(j)}(g))}$ and $\mathsf w_j(1^{(g)})=\mathcal E^{(j)}(\mathfrak M^{(j)}(g))\,g^{-1}$.
\item \emph{(Normalization.)} $\mathfrak m_j(\mathcal E_j(h)g)=h\,\mathfrak m_j(g)$, and $\mathfrak M^{(j)}(\mathsf w_j(U_1))=1$.
\end{enumerate}
\end{theorem}

\begin{proof}
(a), (b) By induction on $j$; the case $j=0$ is trivial. All configurations and transformations below are real, so the transformations preserve the domain of $\ell$, and the relative averages are covariant. Let $U''=U_1^{(\mathsf w_{j-1})}\in\mathrm{Ax}_{j-1}$, so that $U''^{[j-1]}=X_{j-1}=:X$. By Lemma~\ref{lem:axial}(b), the transformations preserving $\mathrm{Ax}_{j-1}$ for $U''$ are $\mathcal E^{(j-1)}(h')$ with $h'$ on $\T^{(j-1)}$, and $(U''^{(\mathcal E^{(j-1)}(h'))})^{[j-1]}=X^{(h')}$. Hence $U_1^{(g)}\in\mathrm{Ax}_j$ iff $g=\mathcal E^{(j-1)}(h')\mathsf w_{j-1}$ with $X^{(h')}$ axial at scale $j$. By Lemma~\ref{lem:axial}(a), the latter $h'$ are exactly $\mathcal E_j(h'')\,g_X$, $h''$ on $\T^{(j)}$, because $X^{(g_X)}$ is axial and the transformations preserving its axiality are the $\mathcal E_j(h'')$. By multiplicativity $g=\mathcal E^{(j)}(h'')\,\mathcal E^{(j-1)}(g_X)\,\mathsf w_{j-1}=\mathcal E^{(j)}(h'')\,\mathsf w_j$. For $g=\mathsf w_j$,
\[
(U_1^{(\mathsf w_j)})^{[j]}=\overline{(U''^{(\mathcal E^{(j-1)}(g_X))})^{[j-1]}}=\overline{X^{(g_X)}}=\tilde{\mathfrak a}_j(X)=\mathcal Q_j(U_1).
\]
(c) $\mathcal E^{(i)}$ does not change values at points of $\T^{(i)}\supset\T^{(j)}$, and $g_{X_i}=\tau_{X_i}^{-1}$ on $\T^{(i+1)}$. (d) If $X$ is axial, then $\tau_X=1$ and $g_X=\mathcal E_j(1)=1$.

(e) One step: let $X=1^{(g)}$ at scale $j-1$. Then $X_W(\Gamma_{y,x})=g(y)R(W(\Gamma_{y,x}))(g(x))^{-1}$. Since $1$ is axial and $X^{(g_X)}=1^{(g_Xg)}$, Lemma~\ref{lem:axial}(a) gives $g_Xg=\mathcal E_j(h)$ with $h=(g_Xg)|_{\T^{(j)}}=\tau_X^{-1}g$. Using $\ell(k^{-1})=-\ell(k)$ and $\ell(RkR^{-1})=R\ell(k)R^{-1}$,
\[
\begin{aligned}
h(y)&=\exp\Big(-\ii L^{-d}\sum_x\ell\big(g(y)R(W(\Gamma_{y,x}))g(x)^{-1}\big)\Big)g(y)\\
&=g(y)\exp\Big(\ii L^{-d}\sum_x\ell\big(g(y)^{-1}R(W(\Gamma_{y,x}))g(x)\big)\Big)=\mathfrak m_j(g)(y).
\end{aligned}
\]
Hence $\tilde{\mathfrak a}_j(1^{(g)})=\overline{1^{(\mathcal E_j(h))}}=1^{(h)}$ and $g_X=\mathcal E_j(\mathfrak m_j(g))g^{-1}$. Iterating, with $X_{j-1}=1^{(\mathfrak M^{(j-1)}(g))}$,
\[
\mathsf w_j(1^{(g)})=\mathcal E^{(j-1)}\big(\mathcal E_j(\mathfrak M^{(j)}g)(\mathfrak M^{(j-1)}g)^{-1}\big)\,\mathcal E^{(j-1)}(\mathfrak M^{(j-1)}g)\,g^{-1}=\mathcal E^{(j)}(\mathfrak M^{(j)}g)\,g^{-1}.
\]
(f) For $g'=\mathcal E_j(h)g$, $g'(y)^{-1}R(W(\Gamma_{y,x}))g'(x)=g(y)^{-1}h(y)^{-1}h(y)R(W(\Gamma_{y,x}))g(x)$, which gives the first identity. For the second, by induction and the first identity, $\mathfrak M^{(j)}(\mathsf w_j)=\mathfrak m_j(g_X\,\mathfrak M^{(j-1)}(\mathsf w_{j-1}))=\mathfrak m_j(g_X)$. By the defining relation of $g_X$, $g_X(y)^{-1}R(W(\Gamma_{y,x}))g_X(x)=X_W(\Gamma_{y,x})$, so $\mathfrak m_j(g_X)=\tau_X^{-1}\tau_X=1$.
\end{proof}

Part~(f) shows that $\mathsf w_k(U_1)$ is the transformation singled out in \cite[(1.29)]{BalReg} by Ba\l aban's normalization $\mathfrak M^{(k)}(g)=1$. We never need that normalization explicitly: the transformation is given by the recursion, and its values at the coarse points by~(c). Parts~(c) and~(d) use only the definitions and Lemma~\ref{lem:axial}(a), so they hold for complex $U_1$ as well, whenever the maps are defined at $U_1$.

\subsection{Expansion of the modified average}

We now expand the modified average around the trivial relative configuration. At scale $j$ write $\delta=\delta_{j-1}$, $\delta'=\delta_j$, $W=W_{j-1}$. For a relative Lie field $B$ at scale $j-1$ and a contour $\gamma=b_1\cdots b_m$ on $\T^{(j-1)}$ let
\[
B^W(\gamma)=\sum_iR(W(\gamma_{<i}))B(b_i),
\]
so that $\ell\big((e^{\ii\delta B})_W(\gamma)\big)=\delta B^W(\gamma)+O((\delta\sum_i|B(b_i)|)^2)$ by Lemma~\ref{lem:log-estimates}(c). Define the \emph{transported tube average}
\[
(q_jB)(c)=L^{-d}\sum_{x\in B_j(c_-)}L^{-1}R\big(W(\Gamma^j_{c_-,x})\big)B^W([x,x(c)])
\]
\cite[(125)]{BalAvg} and the \emph{contour average} $(\gamma_jB)(y)=L^{-d}\sum_{x\in B_j(y)}L^{-1}B^W(\Gamma^j_{y,x})$.

\begin{lemma}[One-step expansion]\label{lem:one-step-expansion}
There are $c_2,C>0$ such that the following holds for $a\le c_2$ and complex $B$ with $|B|\le c_2/\delta'$ on the bonds of $B_j(c_-)\cup B_j(c_+)$ (respectively $B_j(y)$). Put $X=e^{\ii\delta B}$. Then
\[
\tilde{\mathfrak a}_j(X)(c)=e^{\ii\delta'(l_jB+c_j(B))(c)},\qquad \tau_X(y)=e^{\ii\delta'(\gamma_jB+c^\tau_j(B))(y)},
\]
where $l_j=q_j+r_j$ is linear and $c_j,c_j^\tau$ are analytic, with:
\begin{enumerate}[label=(\alph*)]
\item $(l_jB)(c)$ and $c_j(B)(c)$ depend only on $B$ on the bonds of $B_j(c_-)\cup B_j(c_+)$; $\gamma_jB(y)$ and $c^\tau_j(B)(y)$ only on $B$ on the bonds of $B_j(y)$;
\item $|(q_jB)(c)|^2\le L^{-d}\sum_xL^{-1}\sum_{b\in[x,x(c)]}|B(b)|^2$, $\ |(r_jB)(c)|\le Ca\delta'^2\max|B|$, and $|\gamma_jB|\le d\max|B|$;
\item $|c_j(B)|+|c^\tau_j(B)|\le C\delta'\max|B|^2$ and $|c_j(B)-c_j(\tilde B)|\le C\delta'(\max|B|+\max|\tilde B|)\max|B-\tilde B|$ for $\max|B|,\max|\tilde B|\le c_2/(2\delta')$, the maxima being over the relevant bonds;
\item if $W=1$ on the relevant bonds, then $r_j=0$ and $q_j$ is the tube average of Part~I at scale $j$;
\item real $B$ gives $\mathfrak g$-valued $l_jB$, $c_j(B)$, $\gamma_jB$, $c_j^\tau(B)$.
\end{enumerate}
\end{lemma}

\begin{proof}
\emph{Analyticity.} The contours involved have at most $3dL$ bonds, so $\delta\sum|B|\le3dc_2$. Since $W$ is real, $R(W(\cdot))$ is unitary on $\mathfrak g^c$, and all relative transports are products of at most $3dL$ factors $e^{\ii\delta R(W(\cdot))B(b)}$. By Lemma~\ref{lem:log-estimates} and Lemma~\ref{lem:one-step-basic}(a) the arguments of all logarithms in $\tau_X$, $\overline X$ and $\tilde{\mathfrak a}_j(X)$ are within $C(c_2+a\delta'^2)$ of $1$. Hence, for $c_2$ small, these are analytic functions of the finitely many variables $B(b)$, bounded by $C$ in the sense that $|\ell(\tilde{\mathfrak a}_j(X)(c))|\le C(c_2+a\delta'^2)$.

\emph{Remainders.} Let $f(B)=\delta'^{-1}\ell(\tilde{\mathfrak a}_j(e^{\ii\delta B})(c))-(l_jB)(c)$, with $l_j$ the derivative at $B=0$. Then $f(0)=0$ and $df(0)=0$, and $t\mapsto f(tB)$ is analytic for $|t|\le c_2/(\delta'\max|B|)$ with $|f|\le C/\delta'$ there. Cauchy's estimate for the second Taylor coefficient gives $|f(B)|\le C\delta'\max|B|^2$ for $\max|B|\le c_2/(2\delta')$. The Lipschitz bound follows from the Cauchy estimate for $df$. The same argument applies to $\tau_X$.

\emph{The linear part.} Let $X=e^{\ii s\delta B}$ and expand to first order in $s$. Write $h_x=W(\Gamma^j_{c,x})W([c])^{-1}$, $H_x=\ell(h_x)$ and $E=L^{-d}\sum_xH_x$, so that $|H_x|,|E|\le2da\delta'^2$ and $W_j(c)=e^{\ii E}W([c])$. The loops in Definition~\ref{def:one-step} are
\[
(XW)(\Gamma^j_{c,x})(XW)([c])^{-1}=X_W(\Gamma^j_{c,x})\,R(h_x)\big(X_W([c])^{-1}\big)\,h_x .
\]
Let $\mathfrak f(\operatorname{ad}H)$ denote the derivative of $Z\mapsto\ell(e^{\ii Z}e^{\ii H})$ at $Z=0$, and $\mathfrak e(\operatorname{ad}E)$ that of $Z\mapsto\ell(e^{\ii(E+Z)}e^{-\ii E})$; both equal $1$ at $0$. The function $\mathfrak e$ is entire. The function $\mathfrak f$, which is $z\mapsto\ii z/(e^{\ii z}-1)$, is analytic for $|z|<2\pi$ and has poles at the nonzero multiples of $2\pi$; it is used only for $\|\operatorname{ad}H\|\le2|H|$ small. Then $\ell(\overline X(c))=s\delta\big[\mathfrak e(\operatorname{ad}E)\Xi+R(e^{\ii E})B^W([c])\big]+O(s^2)$, with $\Xi=L^{-d}\sum_x\mathfrak f(\operatorname{ad}H_x)\big(B^W(\Gamma^j_{c,x})-R(h_x)B^W([c])\big)$. Moreover $\ell(\tau_X(y))=s\delta L^{-d}\sum_xB^W(\Gamma^j_{y,x})+O(s^2)$. Since $\delta/\delta'=L^{-1}$,
\[
\begin{aligned}
l_jB(c)=L^{-1}\Big\{&\mathfrak e(\operatorname{ad}E)\Xi+R(e^{\ii E})B^W([c])\\
&-L^{-d}\sum_xB^W(\Gamma^j_{c_-,x})+R(W_j(c))L^{-d}\sum_{x'\in B_j(c_+)}B^W(\Gamma^j_{c_+,x'})\Big\}.
\end{aligned}
\]
Decompose
\[
B^W(\Gamma^j_{c,x})=B^W(\Gamma^j_{c_-,x})+R(W(\Gamma^j_{c_-,x}))B^W([x,x(c)])-R(W(\Gamma^j_{c,x}))B^W(\Gamma^j_{c_+,x(c)}).
\] The last term uses $B^W(\gamma^{-1})=-R(W(\gamma))^{-1}B^W(\gamma)$. If all factors $\mathfrak e,\mathfrak f$, $R(h_x)$, $R(e^{\ii E})$ were $1$ and $R(W(\Gamma^j_{c,x}))=R(W_j(c))$, the contour sums would cancel exactly, by the substitution $x'=x(c)$, and only $q_jB$ would remain. The actual differences are controlled by $|\mathfrak f(\operatorname{ad}H_x)-1|\le C|H_x|$, $|\mathfrak e(\operatorname{ad}E)-1|\le C|E|$, $|R(e^{\ii E})-R(h_x)|\le C(|E|+|H_x|)$, and $|R(W(\Gamma^j_{c,x}))-R(W_j(c))|=|R(h_xe^{-\ii E})-1|\le C(|E|+|H_x|)$. Each of them multiplies a contour sum bounded by $3dL\max|B|$. This gives the bound on $r_j=l_j-q_j$. The bound on $q_j$ is the Cauchy--Schwarz inequality, since $R(W(\cdot))$ is unitary, and the bound on $\gamma_j$ is clear. (d) and (e) follow from the formulas.
\end{proof}

\subsection{Expansion of the gauge-adapted averaging}

For a fine field $h$ and a set $S\subset\Tf$ let $\|h\|_S$ be the $\ell^2$ norm over the bonds starting in $S$, as in Part~III. For a bond $c$ of $\T$ put $S_c=B(c_-)\cup B(c_+)$.

\begin{proposition}[Expansion of $\mathcal Q$]\label{prop:Q-expansion}
There are $c_3,C>0$ such that for $a\le c_3$ and complex fine Lie fields $A$ with $|A|_\infty\le c_3$ the following holds. For $j\le k$, $\mathcal Q_j(e^{\ii\eta A})=e^{\ii\delta_jB_j}$ with
\[
B_j=\bQ_jA+C_j(A),\qquad \bQ_j=l_j\cdots l_1,
\]
where $C_j$ is analytic, real for real $A$, and:
\begin{enumerate}[label=(\alph*)]
\item $|B_j|_\infty\le2|A|_\infty$, and $B_j(c)$ depends only on $A$ on the bonds of $B^{(j)}(c_-)\cup B^{(j)}(c_+)$.
\item $|C_j(A)|_\infty\le C|A|_\infty^2$ and $|C_j(A)-C_j(\tilde A)|_\infty\le C(|A|_\infty+|\tilde A|_\infty)|A-\tilde A|_\infty$.
\item \emph{(Kernels.)} For $j=k$ the linear maps $\bQ_k$ and $dC_k(A)$ have kernels, $\bQ_kh(c)=\eta^d\sum_bK(c,b)h(b)$ and $dC_k(A)h(c)=\eta^d\sum_bK_A(c,b)h(b)$, supported on bonds $b$ starting in $S_c$, with
\[
|K(c,b)|\le1+Ca,\qquad|K_A(c,b)|\le C|A|_\infty .
\]
Moreover $|d^2C_k(A)[h,h](c)|\le C\|h\|^2_{S_c}$. Consequently $|\bQ_kh(c)|\le C\|h\|_{S_c}$ and $|dC_k(A)h(c)|\le C|A|_\infty\|h\|_{S_c}$, and the transposes satisfy $|\bQ_k^T\omega|_\infty\le C|\omega|_\infty$ and $|dC_k(A)^T\omega|_\infty\le C|A|_\infty|\omega|_\infty$.
\item The values of the correspondence transformation at the points of $\T$ are $\mathsf w_k(e^{\ii\eta A})(y)=\exp\big(-\ii\,\boldsymbol\Gamma_kA(y)+\ii C^{\mathsf w}_k(A)(y)\big)$, with
\[
\boldsymbol\Gamma_kA=\sum_{i=0}^{k-1}\delta_{i+1}\,\gamma_{i+1}\bQ_iA,\qquad|\boldsymbol\Gamma_kA|\le C|A|_\infty,\qquad|C^{\mathsf w}_k(A)|\le C|A|_\infty^2,
\]
and the analogous Lipschitz bound. At every fine point, $|\ell(\mathsf w_k(e^{\ii\eta A}))|\le C|A|_\infty$.
\end{enumerate}
\end{proposition}

\begin{proof}
(a), (b) By Lemma~\ref{lem:one-step-expansion}, $B_j=l_jB_{j-1}+c_j(B_{j-1})$ with $|l_jB|_\infty\le(1+Ca\delta_j^2)|B|_\infty$ and $|c_j(B)|_\infty\le C\delta_j|B|_\infty^2$. By induction, as long as $|B_{i}|_\infty\le2|A|_\infty$ for $i<j$,
\[
|B_j|_\infty\le|A|_\infty\prod_{i\le j}\big(1+Ca\delta_i^2+2C\delta_i|A|_\infty\big)\le|A|_\infty e^{C(a+|A|_\infty)}\le2|A|_\infty,
\]
using $\sum_i\delta_i\le2$. The hypothesis $|B_{j-1}|\le c_2/\delta_j$ of Lemma~\ref{lem:one-step-expansion} holds since $\delta_j\le1$. Moreover $C_j(A)=\sum_{i\le j}l_j\cdots l_{i+1}c_i(B_{i-1})$, so $|C_j(A)|\le2\sum_iC\delta_i(2|A|)^2\le C|A|^2$. The Lipschitz bound is proved in the same way from the Lipschitz bound in Lemma~\ref{lem:one-step-expansion}(c). Locality follows from Lemma~\ref{lem:one-step-expansion}(a) by induction: the bonds of $B_j(c_-)\cup B_j(c_+)$ have both endpoints in this set, and their scale-$(j-1)$ blocks lie in $B^{(j)}(c_-)\cup B^{(j)}(c_+)$.

(c) For $i\le j$, call a linear map $K$ from scale-$i$ fields to scale-$j$ fields \emph{of type $M$} if $Kv(c')=\delta_i^d\sum_{b'}K(c',b')v(b')$ with $|K(c',b')|\le M\delta_j^{-d}$, and $K(c',b')=0$ unless both endpoints of $b'$ lie in the union of the scale-$j$ blocks of $c'_-$ and $c'_+$. This support is preserved under composition, because the scale-$j$ blocks are unions of scale-$m$ blocks for $m\le j$. At $W=1$ the composition of tube averages from scale $i$ to scale $j$ is the tube average from scale $i$ to scale $j$. By Lemma~\ref{lem:avg}(a), after rescaling, it is of type $1$: its segments $[x,x(c')]$ lie in the two blocks, so both endpoints of each of their bonds do.

\emph{Transports.} For real $W$, $q_j\cdots q_{i+1}$ is a sum over chains of bonds of the $W=1$ weights times products of unitary transports. Hence its kernel is bounded in norm by the $W=1$ kernel, and it is of type $1$.

\emph{Insertions.} Write $l_m=q_m+r_m$. By Lemma~\ref{lem:one-step-expansion}(b), $r_mv(c')$ is a combination of at most $C_L$ values $v(b'')$, with $b''$ a bond of $B_m(c'_-)\cup B_m(c'_+)$ and coefficients bounded by $Ca\delta_m^2$. Expand $l_j\cdots l_{i+1}$ as a sum over the sets $I$ of scales at which $r_m$ is chosen. For $I=\{m\}$ the term is $(q_j\cdots q_{m+1})\,r_m\,(q_{m-1}\cdots q_{i+1})$. In the normalization above, the left factor has kernel at most $\delta_m^d\delta_j^{-d}$ per scale-$m$ bond, and the right factor at most $\delta_i^d\delta_{m-1}^{-d}$ per scale-$i$ bond. For fixed $c',b'$ at most $C_L$ scale-$m$ bonds contribute, so this term is of type $C_La\delta_m^2L^d$. For several insertions, write the term as $A_0r_{m_1}A_1\cdots r_{m_p}A_p$, with $m_1>\dots>m_p$ and transported tube averages $A_q$ of type $1$, and compose from the right. If $T$ is of type $M'$, with the support property above, then $A\,r_m\,T$ is of type $C_La\delta_m^2L^dM'$, by the count for one insertion: for fixed output and input bonds at most $C_L$ (indeed $(2d)^2$) pairs of intermediate bonds of scales $m$ and $m-1$ contribute. Hence the term is of type $\prod_{m\in I}C_La\delta_m^2L^d$. Each counting factor is attached to one insertion; the composition of the factors $A_q$ alone needs none, by the exactness at $W=1$. Altogether $l_j\cdots l_{i+1}$ is of type $\prod_m(1+C_La\delta_m^2)\le1+Ca$. For $i=0$ and $j=k$ this is the bound on $K$.

\emph{First derivatives.} Differentiating the recursion gives $dB_j=(l_j+dc_j(B_{j-1}))\,dB_{j-1}$. By the Cauchy estimate in Lemma~\ref{lem:one-step-expansion}(c), $dc_j(B)v(c')$ is a combination of at most $C_L$ values of $v$, with coefficients bounded by $C\delta_j|B|_\infty\le2C\delta_j|A|_\infty$. So $dc_j$ plays the role of $r_j$, with $a\delta_j^2$ replaced by $C\delta_j|A|_\infty$. The same expansion shows that every product $\prod_m(l_m+dc_m)$ is of type $1+C(a+|A|_\infty)$. The terms with at least one insertion of some $dc_m$ form a map of type $C|A|_\infty$; for $i=0$, $j=k$ they make up $dC_k(A)$. This gives the bound on $K_A$.

\emph{Second derivatives.} The recursion gives
\[
d^2B_j[h,h]=(l_j+dc_j)\,d^2B_{j-1}[h,h]+d^2c_j(B_{j-1})[dB_{j-1}h,dB_{j-1}h].
\]
Since $dB_{j-1}$ is of type $C$ from scale $0$, the Cauchy--Schwarz inequality gives $|dB_{j-1}h(b')|\le C\delta_{j-1}^{-d/2}\|h\|_{S^{(j-1)}_{b'}}$, where $S^{(j-1)}_{b'}=B^{(j-1)}(b'_-)\cup B^{(j-1)}(b'_+)$. By the Cauchy estimate, the source term at a scale-$j$ bond $c'$ is therefore bounded by $C_L\delta_j\delta_{j-1}^{-d}\|h\|^2_{S^{(j)}_{c'}}$. It is propagated to $\T$ by a map of type $C$ from scale $j$, so its contribution at $c$ is at most
\[
C\sum_{c'}\delta_j^d\,C_L\delta_j\delta_{j-1}^{-d}\|h\|^2_{S^{(j)}_{c'}}\le C_L\delta_j\|h\|^2_{S_c},
\]
because each fine bond lies in a bounded number of the sets $S^{(j)}_{c'}$, all of which lie in $S_c$. Summing over $j$ gives the bound on $d^2C_k$.

The remaining consequences follow from the kernel bounds. The bounds on $\bQ_kh(c)$ and $dC_k(A)h(c)$ follow by the Cauchy--Schwarz inequality, since $S_c$ has volume $2$. The bounds on the transposes follow by summing the kernel bounds over $c$, since each fine bond lies in at most $2d$ sets $S_c$.

(d) By Theorem~\ref{thm:correspondence}(c) and Lemma~\ref{lem:one-step-expansion}, $\ell(\tau_{X_i})=\delta_{i+1}(\gamma_{i+1}B_i+c^\tau_{i+1}(B_i))$, of size at most $C\delta_{i+1}|A|$. Lemma~\ref{lem:log-estimates}(c) applied to the product gives the formula and the bounds, since $\sum_i\delta_{i+1}\le2$. At a fine point $x$, $\mathsf w_k=\prod_i\mathcal E^{(i)}(g_{X_i})$, and $g_{X_i}(x)=R(W_i(\Gamma))^{-1}(\tau_{X_i}(y)^{-1}(X_i)_W(\Gamma))$ for $x\in B_{i+1}(y)$ has logarithm at most $C\delta_{i+1}|A|$. Conjugation by $R(W)$ does not change the norm.
\end{proof}

\subsection{The linearized averages}

Let $q'_j\lambda(y)=L^{-d}\sum_{x\in B_j(y)}R(W_{j-1}(\Gamma^j_{y,x}))\lambda(x)$ be the transported block average of scalar fields, and put
\[
Q'_k(U_0)=q'_k\cdots q'_1,\qquad\bQ_k(U_0)=l_k\cdots l_1 .
\]
Explicitly, $(Q'_k\lambda)(y)=\eta^d\sum_{x\in B(y)}T(y,x)\lambda(x)$, where
\[
T(y,x)=R(W_{k-1}(\Gamma^k_{y,x_{k-1}}))\cdots R(W_0(\Gamma^1_{x_1,x}))
\]
is the transport along the multiscale tree, with the averaged backgrounds at each scale. In particular $Q'^{\,T}_k=\mathcal E^{(k)}$. These are Ba\l aban's linearized averages \cite[(3.15)--(3.19)]{BalBackground}.

\begin{proposition}[Linearized averages]\label{prop:linearized-averages}
Let $U_0\in\mathfrak U_0(a)$ be real with $a\le c_3$.
\begin{enumerate}[label=(\alph*)]
\item $Q'_kQ'^{\,T}_k=I$.
\item \emph{(Exact covariance.)} $\bQ_k(U_0)D_{U_0}=D_{W_k}Q'_k(U_0)$. In the notation of Part~II, the defect vanishes: $E=0$, with $V=W_k$.
\item At $U_0=1$: $Q'_k=Q_0$ and $\bQ_k=\bQ$, the operators of Part~I.
\item \emph{(Locality.)} $(\bQ_kA)(c)$ depends only on $A$ and $U_0$ on the bonds of $B(c_-)\cup B(c_+)$, and $(Q'_k\lambda)(y)$ only on $\lambda$ and $U_0$ on $B(y)$.
\item \emph{(Closeness.)} Suppose that on a cube containing $B(c_-)\cup B(c_+)$ we have $U_0=e^{\ii\eta A_0}$ (in some real gauge) with $|A_0|\le\eps$, and $a\le\eps$, where $\eps\le c_3$. Then, computed in this gauge, the kernels satisfy $|K_{\bQ_k(U_0)}(c,b)-K_{\bQ}(c,b)|\le C\eps$ and $|K_{Q'_k(U_0)}(y,x)-K_{Q_0}(y,x)|\le C\eps$. Hence the differences are bounded by $C\eps$ from sup norms to sup norms and in the block norms of Lemma~\ref{lem:cov-avg}(b), and so are their transposes.
\item \emph{(Covariance.)} For every real gauge transformation $u$ on $\Tf$: $W_k(U_0^u)=W_k(U_0)^{u|_\T}$, $Q'_k(U_0^u)R(u)=R(u|_\T)Q'_k(U_0)$ and $\bQ_k(U_0^u)R(u)=R(u|_\T)\bQ_k(U_0)$.
\end{enumerate}
\end{proposition}

\begin{proof}
(a) $q'^{\,T}_jg(x)=R(W_{j-1}(\Gamma^j_{y,x}))^{-1}g(y)$, because the volume factors $L^{-d}\delta_j^d/\delta_{j-1}^d$ equal $1$. Hence $q'_jq'^{\,T}_j=I$, since $R(W)$ is unitary.

(b) Let $\lambda$ be real and $g_t=e^{\ii t\lambda}$. Then $1^{(g_t)}=e^{-\ii t\eta D_{U_0}\lambda+O(t^2)}$ on $\Tf$. By Theorem~\ref{thm:correspondence}(e), $\mathcal Q_k(1^{(g_t)})=1^{(\mathfrak M^{(k)}(g_t))}$ on $\T$. Differentiating the definition of $\mathfrak m_j$ at $g=1$ gives $\mathfrak M^{(k)}(g_t)=e^{\ii tQ'_k\lambda+O(t^2)}$, hence $1^{(\mathfrak M^{(k)}(g_t))}=e^{-\ii tD_{W_k}Q'_k\lambda+O(t^2)}$. Comparing the terms of first order in $t$ gives $\bQ_kD_{U_0}\lambda=D_{W_k}Q'_k\lambda$.

(c) At $U_0=1$, $q'_j$ is the block average at scale $j$, and $l_j=q_j$ is the tube average at scale $j$ by Lemma~\ref{lem:one-step-expansion}(d). Compositions of block averages are block averages. Compositions of tube averages are tube averages: in $(q_2q_1A)(y)$ the shifts $m\delta_1+m'\delta_0$, $0\le m,m'<L$, run through $0,\delta_0,\dots,(L^2-1)\delta_0$ exactly once. By induction, $\bQ_k$ is the tube average $\bQ$ of Definition~\ref{def:averages}.

(d) follows from Lemma~\ref{lem:one-step-expansion}(a) and Lemma~\ref{lem:one-step-basic}(c).

(e) By (f) we may compute in the given gauge. By Lemma~\ref{lem:avg-local-gauge}, $W_{j}=e^{\ii\delta_jA_{0,j}}$ on the relevant bonds with $|A_{0,j}|\le C\eps$. Hence all transports $R(W_{j-1}(\gamma))$ along contours of at most $3dL$ bonds differ from $1$ by at most $C\eps\delta_j$. Together with the bound on $r_j$ in Lemma~\ref{lem:one-step-expansion}(b), this shows that $l_j-q_j(1)$ is a combination of at most $C_L$ values with coefficients $C\eps\delta_j$. Telescope:
\[
\bQ_k(U_0)-\bQ=\sum_jl_k\cdots l_{j+1}\big(l_j-q_j(1)\big)q_{j-1}(1)\cdots q_1(1).
\]
By the bookkeeping of Proposition~\ref{prop:Q-expansion}(c), the kernel of the difference is bounded by $\sum_jC_L\eps\delta_j\le C\eps$. The same argument applies to $Q'_k$. Kernel bounds give the sup-norm bounds, the block norms (by the Schur test), and the transposes.

(f) follows from Lemma~\ref{lem:one-step-basic}(b), the definitions, and the transformation law of the relative gauge action.
\end{proof}

\begin{corollary}[Parts~II and~III apply to Ba\l aban's averaging]\label{cor:parts-apply}
Let $U_0$ be real and $(M,\eps)$-regular with $\eps$ small. Replace the averaging pair $(Q'(U),\bQ(U))$ by $(Q'_k(U_0),\bQ_k(U_0))$, the defect $E$ by $0$, and the unit-lattice background $V$ by $W_k$. Then the real results of Sections~\ref{sec:cov-scalar}--\ref{sec:reg-gauge} hold, and so does Theorem~\ref{thm:random-walk}. Proposition~\ref{prop:reg-complex} holds for $\delta\le\min(\delta_0,\delta_{AB}/C)$ with $\eps+C\delta\le\min(\eps_0,c'_3)$; at the perturbations $U_\delta$ the averages are those of Proposition~\ref{prop:AB}.
\end{corollary}

\begin{proof}
An $(M,\eps)$-regular background has curvature at most $C\eps$, by Lemma~\ref{lem:curvature} in the local gauges. Proposition~\ref{prop:linearized-averages} gives the properties (a)--(d) of Lemma~\ref{lem:cov-avg}, in each local gauge, with $\eps$ replaced by $C\eps$. By Remark~\ref{rem:general-averaging}, these are the only properties of the averaging used in the real results of Parts~II and~III. Property (v) of that remark at the $U_\delta$ is shown there, from Proposition~\ref{prop:AB}(iv)--(vi) and the identity theorem. Proposition~\ref{prop:AB} uses Part~IV only at the real reference, so there is no circularity. It controls only the sup norm of the perturbation; the derivative and H\"older bounds on $\delta\!A$ are hypotheses of Proposition~\ref{prop:reg-complex}.
\end{proof}

From now on $Q'=Q'_k(U_0)$, $\bQ=\bQ_k(U_0)$, and all operators of Parts~II and~III ($G'$, $R$, $H'$, $\Gamma'$, $\bG$, $\bH$, $\bC$, \dots) are taken at the background $U_0$ with this averaging pair. Since $E=0$, the slice projection satisfies $\bQ\Pi=\bQ$ (Lemma~\ref{lem:slice}(c)), and the correction $\mathcal K_2$ of Section~\ref{sec:cov-gauge} vanishes.
